\documentclass[10pt,journal,onecolumn]{IEEEtran}

% ===================== 宏包 =====================
\usepackage{amsmath,amssymb,amsfonts,amsthm}
\usepackage{algorithm,algorithmic}
\usepackage{graphicx}
\usepackage{booktabs,multirow,array}
\usepackage[table]{xcolor}
\usepackage{colortbl}
\usepackage{pifont}
\usepackage{cite}
\usepackage{url}
\usepackage[hidelinks]{hyperref}
\usepackage{amsmath, amssymb}

\newtheorem{definition}{Definition}
\newtheorem{proposition}{Proposition}
\newtheorem{lemma}{Lemma}
\newtheorem{assumption}{Assumption}

\definecolor{best}{HTML}{D4EDDA}
\definecolor{second}{HTML}{D1ECF1}
\definecolor{rowhi}{HTML}{FFF3CD}
\definecolor{negbad}{HTML}{F8D7DA}
\definecolor{posgood}{HTML}{D4EDDA}
\definecolor{tabgray}{HTML}{F5F5F5}

\newcommand{\cmark}{\checkmark}
\newcommand{\xmark}{$\times$}
\newcommand{\best}[1]{\cellcolor{best}{\bfseries #1}}
\newcommand{\secondb}[1]{\cellcolor{second}#1}
\newcommand{\pos}[1]{\textcolor{green!50!black}{{\bfseries #1}}}
\newcommand{\bad}[1]{\textcolor{red!70!black}{#1}}
\newcommand{\dON}[1]{\cellcolor{green!18}{\bfseries #1}}
\newcommand{\dOFF}[1]{\cellcolor{gray!22}#1}
\newcommand{\gcheck}{{\color{green!50!black}\ding{51}}}
\newcommand{\rcross}{{\color{gray!65}\ding{55}}}
\newcommand{\heat}[2]{\cellcolor{blue!#1!white}#2}
\newcommand{\ind}{\mathbb{I}}

\newcommand{\system}{SpecSched}

\begin{document}

\title{Supplementary Appendix for \\ SpecSched: Cache-Coupled Speculative Scheduling for Multi-Tenant LLM Serving}
\author{}
\maketitle
\appendices

\section{Stress Coverage, Overhead, and Artifact Scope}
\label{app:stress-overhead-artifact}

\subsection{Stress Coverage and Failure Modes}
\label{sec:stress-failure}

The stress evidence separates three collapse modes rather than treating
all high-load failures as one phenomenon. First, verifier saturation
appears when long drafts cannot be converted into efficient verification
batches. The NoBatch ablation in the main-paper ablation results reduces
goodput by about one third at both $N=32$ and $N=64$, showing that
accepted-token potential is insufficient without batching control.
Second, KV thrashing appears when cache residency is invisible to the
speculative scheduler. The NoEvict row at $N=64$ keeps raw progress
possible but loses 68.7\% SLO-aware goodput and 0.64 SLO-attainment
points because requests wait behind a saturated cache. Third,
speculation collapse appears when the draft path itself is
unprofitable. The warm-up detector disables unsafe local model pairs and
the Jetson remote path before they consume SLO slack. These measured
cases support the local control claim: \system{} needs batching,
resident-aware cache control, and fallback detection together.

The route variable adds a distributed failure surface. A route-aware
deployment can add two additional modes: routing oscillation, where
requests move between verifier pools faster than KV reuse can amortize
migration, and remote-KV degradation, where a higher remote hit rate is
offset by fetch latency. The current artifact cannot execute those
cluster-level controls natively, so Appendix~\ref{app:route-kv-sim}
reports a calibrated route/KV stress simulation. It supports the
mechanism that route cost and cache residency interact, but it is not
used as a native distributed-serving performance claim.

\subsection{Scheduling Overhead}
\label{sec:overhead}

The implementation records CPU-side timing for candidate selection,
marginal allocation, KV accounting, and cleanup. A dedicated
microbenchmark sweeps $N\in\{8,16,32,64\}$ and
$k_{\max}\in\{3,5,7\}$ on the local 0.5B/0.5B pair. Mean scheduler
overhead is 123--295~$\mu$s per epoch, and the end-to-end scheduler
fraction stays below 0.3\% in every cell, consistent with the
$O(N k_{\max}+N\log N)$ bound proved in the main text.
We use this microbenchmark as a reproducibility and scaling check, not
a headline performance result.

\subsection{Artifact Scope}
\label{sec:artifact-scope}

The artifact supports result-level reproducibility. It includes
experiment configurations, saved raw and processed CSV/JSON summaries,
plotting scripts, generated figures and tables, and claim-audit checks.
From saved outputs, \texttt{make tpds-results-reproduce} regenerates
the reported plots and tables and verifies that headline numerical
claims match the processed results. The Docker path is intentionally
saved-output only: it runs without CUDA, PyTorch, vLLM, model
checkpoints, traces, or Jetson access. Full GPU and edge reruns require
the listed hardware, model checkpoints, trace inputs, and
CUDA/PyTorch/vLLM environment. Dockerized one-command GPU reruns and
formal artifact-evaluation packaging remain outside the submitted
evidence.

\section{Evaluation Matrix and Setup Details}
\label{app:evaluation-scope}

This appendix records scope, notation, and diagnostic evidence that
would interrupt the main experimental narrative. It does not add new
headline claims beyond the measured vLLM-based artifact. In particular,
the current manuscript does not include production-trace generality,
native speculative-serving comparisons, multi-GPU tensor parallelism,
optimized edge RPC, 24-hour cost or energy studies, or Dockerized
one-command GPU reruns. The model-pair, pressure, overhead, and
open-loop controller tables below are therefore boundary diagnostics
and reproducibility aids, not additional leaderboards.

The current artifact does not include production traces such as Azure
LLM inference traces, BurstGPT, or WildChat, nor does it include
2-GPU tensor-parallel verification, cross-host profitability beyond
the Jetson SSH/JSONL path, or 24-hour endurance and cost/energy
studies. TETRIS, EAGLE-2, and SpecServe are common-executor scheduler
proxies. They do not replace native integrations with feature-level
draft generators, tree drafting, optimized communication, or runtime
overlap. Native EAGLE-2/SpecInfer reproduction, native AdaSpec
integration, composition with SGLang/RadixAttention and LMCache-style
reusable KV layers, disaggregated prefill/decode integration,
Dockerized one-command reruns, GPU reruns, and formal
artifact-evaluation packaging remain engineering extensions rather
than current experimental evidence. The result-level artifact checks
saved outputs and claims but does not convert those missing
experiments into evidence.

The same scope boundary applies to system-wide KV-cache control. The
route/KV simulation adds a controlled stress surface over route delay,
remote-hit ratio, migration, concurrency, and KV budget, but it is not
a native cluster measurement. The current artifact still does not
measure real remote cache hits, inter-GPU KV migration, request
routing among verifier pools, or prefill/decode tier placement. Those
native dimensions remain extension work before claiming
distributed-control stability. We also use several 2026 arXiv
preprints as related-work signposts for fast-moving speculative
serving and concurrency-control work. They were verified as preprints
during manuscript preparation, but they should not be read as archival
baselines. Native-system comparisons against those systems remain out
of scope for the current artifact.

Table~\ref{tab:evidence-roadmap} maps the main reviewer-risk items to
their current evidence status. Its purpose is to make the paper's scope
auditable: rows marked as measured support claims in the main text,
whereas rows marked as future work require new experiments before they
can support stronger conclusions.

\subsection{Single-User Draft-Length Reference}
\label{app:single-user-opt}

For a single user $i$ with acceptance rate $\alpha_i$ and exclusive
access to the target model, the standard draft-length objective is
\begin{equation}
    \label{eq:single-opt}
    k_i^* = \arg\max_{k} \;
    \frac{G_i(k)}{c_d \cdot k + c_v(k)},
\end{equation}
where $c_d$ is the per-token draft cost and $c_v(k)$ is the
verification cost for $k$ tokens. Under the i.i.d.\ approximation
($\alpha_\ell=\alpha_i$ for all $\ell$), the familiar closed-form
rule~\cite{leviathan2023fast,xia2024unlocking} is
\begin{equation}
    \label{eq:single-opt-closed}
    k_i^* =
    \left\lfloor
    \frac{-1}{\ln \alpha_i}
    \cdot \ln\!\left(\frac{c_v'(0)}{c_d}\right)
    \right\rfloor .
\end{equation}
This reference is useful only as a contrast: in shared serving,
$c_v(k)$ is not private to request $i$, because batch size, padded
draft length, and the other active users jointly determine the
realized verifier cost.

\begin{table*}[t]
\centering
\caption{Evidence roadmap for the revised TPDS framing. Measured rows
are supported by the current artifact. Future-work rows are intentionally
not used as evidence for headline claims.}
\label{tab:evidence-roadmap}
\scriptsize
\setlength{\tabcolsep}{3.0pt}
\renewcommand{\arraystretch}{1.12}
\begin{tabular}{@{}p{0.20\textwidth}p{0.34\textwidth}p{0.18\textwidth}p{0.20\textwidth}@{}}
\toprule
\textit{Reviewer-risk item} & \textit{Current manuscript handling} &
\textit{Status} & \textit{Claim supported} \\
\midrule
BestFixed comparison & Three-seed closed-loop and Poisson paired CIs
cross zero against the offline fixed-$k$ envelope & measured boundary &
Online control approaches the tuned envelope, not dominance \\
BestFixed transfer stress & Calibrated diagnostic over workload shift,
bursty arrivals, and tiered SLOs compares one global fixed depth,
phase-tuned fixed depth, and \system{} & simulation diagnostic &
Online adaptation advantage when a single fixed depth is mismatched \\
Measured transfer replay & Three-seed GPU replay covers workload shift,
bursty arrivals, and tiered SLOs at $N=64$ against the best full-run
fixed depth from $k\in\{1,3,5,7\}$. A tiered-SLO supplement also
includes \system{}-NoSLO & measured diagnostic, not headline evidence &
Current measured replay is mixed-to-negative, so measured replay is the
authoritative boundary for the transfer claim \\
Proxy baseline scope & TETRIS, EAGLE-2, and SpecServe rows use the same
executor, traces, KV manager, and SLO calibration & measured proxy &
Scheduler-signal comparison only \\
KV eviction ranking & Three-seed KV-pressure cells compare LRU,
longest-context-first, H2O-style, vLLM-style, and future-value eviction
under selected budgets & measured boundary & Cache residency should be
visible, but no static rule is universal \\
Unsafe model pairs and edge path & Warm-up detector results cover unsafe
Qwen2.5-1.5B/0.5B, SmolLM2, profitable Pythia, and Jetson remote drafting
& measured safety & Disable unprofitable speculation before it consumes
SLO slack \\
Scheduler overhead & CPU-side microbenchmark sweeps $N$ and $k_{\max}$
on a local model pair & measured diagnostic & Online bookkeeping is not
the observed bottleneck \\
Routing and remote-KV stress & Calibrated simulation sweeps
local-only, random, KV-aware, and oracle routes under route delay, KV
budget, and concurrency & simulation diagnostic & Route/KV/load coupling
and oscillation failure modes \\
\midrule
Fine-grained component ablation & Would separate resident weighting,
re-prefill budget, future-value eviction, padding-aware batching, and
detector effects individually & future work & Per-module causal
attribution \\
Native or faithful speculative baselines & Would integrate native
TETRIS, EAGLE-2/SpecInfer, or AdaSpec/SpecServe runtime behavior &
future work & End-to-end system comparison \\
Datacenter or multi-GPU serving & Would repeat target/draft and
prefill/decode disaggregation on A100/H100/L40S-class or multi-GPU pools
& future work & Hardware-general performance claim \\
Phase-tuned transfer oracle on GPU & Would run per-phase fixed-depth
oracles, recovery-time metrics, and paired confidence intervals across
the three measured non-stationary traces & future work &
Stronger measured separation between one global fixed depth and online
adaptation \\
Native routing-aware speculative scheduling & Would add real model or
verifier-pool routing, measured remote-KV hits, and routing stability
constraints & future work & Native distributed serving control claim \\
Native disaggregated KV-cache placement & Would measure remote cache
fetch latency, migration cost, and prefill/decode tier placement in a
cluster runtime & future work & System-wide KV-resource claim \\
Production joint stress & Would add burst intensity and real cluster
routing to the current simulation surface & future work & Stability
under coupled production stress \\
\bottomrule
\end{tabular}
\end{table*}

\begin{figure*}[t]
    \centering
    \includegraphics[width=\textwidth]{figures_scientific/bestfixed_stress_scientific.pdf}
    \caption{Calibrated BestFixed transfer stress diagnostic at $N=64$.
    The simulation contrasts one global fixed draft depth with a
    phase-tuned fixed-depth oracle and \system{} under workload shift,
    bursty arrivals, and tiered SLOs. It is a mechanism illustration,
    not headline measured evidence.}
    \label{fig:bestfixed-stress}
\end{figure*}

% Inlined from ../paper_assets/tables/bestfixed_stress_sim.tex.
\begin{table}[t]
\centering
\caption{Calibrated BestFixed transfer stress diagnostic at $N=64$.
Values are simulation means over five deterministic replicates.}
\label{tab:bestfixed-stress}
\scriptsize
\setlength{\tabcolsep}{3.0pt}
\renewcommand{\arraystretch}{1.12}
\begin{tabular}{@{}llrrr@{}}
\toprule
Scenario & Method & Goodput $\uparrow$ & P99 $\downarrow$ & Recovery $\downarrow$ \\
\midrule
Shift & BestFixed & 66.0 & 1247 & 45 \\
Shift & \system{} & 78.1 & 1289 & 18 \\
Bursty & BestFixed & 67.8 & 1256 & 25 \\
Bursty & \system{} & 78.3 & 1250 & 10 \\
Tiered & BestFixed & 69.1 & 1245 & 0 \\
Tiered & \system{} & 80.1 & 1245 & 0 \\
\bottomrule
\end{tabular}
\vspace{-0.5em}
\end{table}

Table~\ref{tab:param-summary} summarizes the symbols and defaults used
by the formulation, algorithm, and experiments. Unless otherwise
stated, the scheduler uses $\gamma=0.3$, warm-starts
$\hat{\alpha}_{i}=0.8$, and sets $\lambda_v=1$, $\eta=0.2$,
$\zeta=0.1$, $\lambda_p=0.1$, and $\lambda_e=0.5$. Runtime quantities
such as $N$, $C_v(t)$, $M_{\text{KV}}$, $Q_i(t)$, and
\textsc{KeepValue} are measured or computed per experiment rather than
tuned per baseline.

\begin{table*}[t]
\centering
\caption{Parameter and symbol summary. Defaults apply unless a
specific experiment overrides them. Runtime state variables are logged
per run rather than fixed globally.}
\label{tab:param-summary}
\scriptsize
\setlength{\tabcolsep}{3.0pt}
\renewcommand{\arraystretch}{1.12}
\begin{tabular}{@{}p{0.13\textwidth}p{0.32\textwidth}p{0.18\textwidth}p{0.27\textwidth}@{}}
\toprule
\textit{Symbol} & \textit{Meaning} & \textit{Location} & \textit{Default or role} \\
\midrule
\multicolumn{4}{@{}l}{\textit{Workload and SLO variables}} \\
\cmidrule(l){1-4}
$N$ & Number of concurrent users or active requests & Evaluation & Swept over $4,8,16,32,64$. Main three-seed claims use $16,32,64$ \\
$\tau_i^{\text{SLO}}$ & TPOT target for user $i$ & Main-text SLO definition & Calibrated from NoSpec by output-length bucket \\
$\beta_i$ & SLO multiplier used to set $\tau_i^{\text{SLO}}$ & Main-text SLO definition and setup & Medium setting $1.5$. Sensitivity uses $\{1.25,1.5,2.0\}$ \\
$Q_i(t)$ & Virtual queue measuring accumulated SLO risk & Main-text virtual queue & Runtime state, capped to prevent monopolization \\
$L_i(t)$ & Context length of request $i$ & System model & Runtime state updated after accepted tokens \\
$r_i(t)$ & Serving route for request $i$ & Main-text route capacity and joint objective & Fixed to $r_{\mathrm{local}}$ in the measured artifact \\
\midrule
\multicolumn{4}{@{}l}{\textit{Speculation and verifier controls}} \\
\cmidrule(l){1-4}
$k_i$ & Draft length assigned to request $i$ & Main-text joint objective & Chosen online from $\{0,\ldots,k_{\max}\}$ \\
$k_{\max}$ & Maximum draft length considered by the scheduler & Algorithm and overhead study & Configuration cap. Overhead study sweeps $3,5,7$ \\
$\hat{\alpha}_{i}(t)$ & Estimated draft-token acceptance probability & Main-text acceptance estimator & Warm-start $0.8$, updated by EMA \\
$\gamma$ & EMA smoothing parameter for acceptance estimates & Main-text acceptance estimator & $0.3$ \\
$C_v(t)$ & Verifier-token capacity available in epoch $t$ & Resource constraint & Profiled runtime budget \\
$C_v^r(t)$ & Route-specific verifier-token capacity & Main-text compute constraint & Reduces to $C_v(t)$ for $r_{\mathrm{local}}$ \\
$\lambda_v$ & Weight on verifier-load cost in the objective & Main-text joint objective & $1$ \\
$\lambda_p$ & Weight on padding cost in the objective & Main-text joint objective & $0.1$ \\
$\ell_r(t)$, $\lambda_r$ & Route cost and its objective weight & Main-text eviction cost and joint objective & $\ell_r=0$ in current experiments, nonzero for remote KV or migration studies \\
\midrule
\multicolumn{4}{@{}l}{\textit{KV-cache and fallback controls}} \\
\cmidrule(l){1-4}
$M_{\text{KV}}$ & KV-cache block budget & Main-text memory constraint & Configured per experiment \\
$M_{\text{perm}}$, $M_{\text{temp}}$, $M_{\text{act}}$ & Permanent KV, temporary speculative KV, and activation/scratch usage & Main-text KV decomposition & Runtime memory accounting \\
$h_i(t)$ & Indicator that request $i$'s KV-cache is resident & System state & $1$ resident, $0$ evicted \\
$R_{\text{prefill}}(t)$ & Re-prefill admission budget per epoch & Main-text re-prefill budget & Derived from $r_{\max}C_v(t)$ \\
$\eta$ & Weight on incremental KV pressure in marginal allocation & Main-text marginal score & $0.2$ \\
$\zeta$ & Weight on context-length imbalance in batch composition & Main-text padding cost & $0.1$ \\
\textsc{KeepValue}$_i(t)$ & Expected future generated-token value per resident KV and route cost & Main-text eviction cost & Computed online from $\bar{w}_iG_i(k_i^\star,t)/(b_i(L_i)+\mu\ell_{r_i})$ \\
$\lambda_e$ & Weight on resident cache value in the objective & Main-text joint objective & $0.5$ \\
$\epsilon$ & Minimum warm-up throughput margin for enabling speculation & Main-text speculation detector & $0.05$ \\
$\rho_{\text{prefill}}(t)$, $\theta_{\text{prefill}}$ & Sliding-window re-prefill rate and overload threshold & Main-text re-prefill monitor & Diagnostic degradation hook, not a headline-tuned result \\
\bottomrule
\end{tabular}
\vspace{-0.5em}
\end{table*}

Table~\ref{tab:exp-setup} summarizes the measured artifact coverage.
The main scheduler and KV-pressure results use the vLLM-based
scheduler path. Each run records seed, model path, SLO mode, KV
budget, maximum batch size, peak GPU allocation, per-epoch scheduler
overhead, and cache events. The logged verifier-load ratio is a
pre-trimming pressure indicator, not a claim about executed
target-batch infeasibility.

\begin{table*}[t]
\centering
\caption{Evaluation groups and evidence scope in the current artifact.
Three-seed groups support the main statistical claims. Diagnostic
groups bound practical failure modes and are not presented as
production-trace or native-system comparisons.}
\label{tab:exp-setup}
\scriptsize
\setlength{\tabcolsep}{2.0pt}
\renewcommand{\arraystretch}{1.12}
\begin{tabular}{@{}p{0.10\textwidth}p{0.13\textwidth}p{0.17\textwidth}p{0.16\textwidth}p{0.07\textwidth}p{0.30\textwidth}@{}}
\toprule
\textit{Role} & \textit{Hardware} & \textit{Target / draft} &
\textit{Workload} & \textit{$N$} & \textit{Claim / scope} \\
\midrule
\multicolumn{6}{@{}l}{\textit{Headline comparisons}} \\
\cmidrule(l){1-6}
\rowcolor{rowhi}Scheduler & RTX~3090~Ti & Qwen2.5-7B / Qwen2.5-0.5B &
ShareGPT/LMSYS-style output lengths &
4, 8, 16, 32, 64 & Common-executor scheduler comparison under calibrated SLOs \\
\rowcolor{rowhi}KV pressure & RTX~3090~Ti & Qwen2.5-7B / Qwen2.5-0.5B &
Long-output ShareGPT/LMSYS traces &
16, 32, 64 & Cache-regime boundary with no universal eviction-rule claim \\
\midrule
\multicolumn{6}{@{}l}{\textit{Boundary and safety diagnostics}} \\
\cmidrule(l){1-6}
Model pairs & RTX~3090~Ti &
Qwen2.5-1.5B/0.5B, SmolLM2-1.7B/135M, Pythia-2.8B/160M &
Short ShareGPT/LMSYS traces &
16, 32 & Model-pair safety and transfer boundaries \\
Edge practical & RTX~3090~Ti + 3 Jetsons & Qwen2.5-7B / Qwen2.5-0.5B &
ShareGPT, Tailscale/SSH draft RPC &
4, 8 & Safe fallback on a real remote-draft path \\
\midrule
\multicolumn{6}{@{}l}{\textit{Illustration and cost checks}} \\
\cmidrule(l){1-6}
Motivation & RTX~3090~Ti + Jetson & Qwen2.5-7B / Qwen2.5-0.5B &
ShareGPT, 128-token prompt cap, 8 output tokens &
4, 8, 16, 32 & Edge overhead diagnostic, not an optimized edge-serving result \\
Overhead & RTX~3090~Ti & Qwen2.5-0.5B / Qwen2.5-0.5B &
Synthetic scheduler microbenchmark &
8, 16, 32, 64 & CPU scheduler-cost scaling only \\
\bottomrule
\end{tabular}
\end{table*}

The appendix tables should be read as boundary documentation rather
than as additional headline leaderboards. Table~\ref{tab:exp-setup}
separates evidence used for the main scheduler claims from diagnostics
that explain when the same control logic falls back, when cache pressure
changes the preferred eviction rule, and how much CPU-side scheduling
cost is added by the policy. This separation keeps the main text focused
on calibrated three-seed comparisons while preserving the setup details
needed to audit model-pair transfer, pressure sensitivity, and overhead.

Tables~\ref{tab:model-pair-performance} and
\ref{tab:model-pair-n32} report model-pair safety checks, so each row
should be interpreted within its target/draft pair rather than as a
cross-family ranking. At $N=16$, the detector disables speculation for
the Qwen2.5-1.5B/0.5B and SmolLM2-1.7B/135M pairs, where the safe path
stays close to NoSpec instead of following harmful speculative settings.
For Pythia-2.8B/160M, the detector keeps speculation enabled and
improves over NoSpec without claiming superiority over the offline-tuned
BestFixed envelope or the TETRIS proxy. The $N=32$ table records the
same boundary at higher concurrency: the Qwen2.5 pair remains unsafe for
speculation, whereas the Pythia pair remains profitable.

% Inlined from ../paper_assets/tables/model_pair_performance.tex.
\begin{table*}[t]
\centering
\caption{Model-pair sensitivity and safe fallback at $N=16$ across three seeds. Values are SLO-aware goodput in tok/s ($\uparrow$), mean$\pm$std. Detector color gives the warm-up decision. The safe-fallback column reports paired lift against NoSpec.}
\label{tab:model-pair-performance}
\scriptsize
\setlength{\tabcolsep}{2.5pt}
\renewcommand{\arraystretch}{1.12}
\begin{tabular}{@{}p{0.10\textwidth}p{0.15\textwidth}ccccc cc@{}}
\toprule
\multirow{2}{*}{Family} & \multirow{2}{*}{Target / draft} &
\multicolumn{3}{c}{Baselines ($\uparrow$)} &
\multicolumn{2}{c}{\system{} variants ($\uparrow$)} &
\multicolumn{2}{c}{Safety outcome} \\
\cmidrule(lr){3-5}\cmidrule(lr){6-7}\cmidrule(l){8-9}
 & & NoSpec & BestFixed & TETRIS-Proxy & plain & safe & Detector & Safe vs NoSpec \\
\midrule
Qwen2.5 & 1.5B / 0.5B & {\bfseries 99.4$\pm$5.9} & 80.3$\pm$7.4 & 65.6$\pm$5.1 & 70.0$\pm$8.8 & 98.8$\pm$6.3 & \dOFF{OFF} & \secondb{$-$0.7\%} \\
SmolLM2~\cite{allal2025smollm2} & 1.7B / 135M & {\bfseries 120.5$\pm$8.4} & 41.5$\pm$3.4 & 33.0$\pm$6.2 & 23.8$\pm$13.5 & 115.0$\pm$13.7 & \dOFF{OFF} & \secondb{$-$4.7\%} \\
Pythia~\cite{biderman2023pythia} & 2.8B / 160M & 87.0$\pm$1.2 & 131.4$\pm$7.0 & {\bfseries 134.4$\pm$17.6} & 126.1$\pm$11.4 & 127.5$\pm$15.0 & \dON{ON} & \pos{+46.6\%} \\
\bottomrule
\end{tabular}
\vspace{-0.5em}
\end{table*}

% Inlined from ../paper_assets/tables/model_pair_n32.tex.
\begin{table}[t]
\centering
\caption{Additional $N=32$ model-pair boundary checks across three seeds. Values are SLO-aware goodput in tok/s ($\uparrow$), mean$\pm$std. The last column reports paired \system{} lift against NoSpec.}
\label{tab:model-pair-n32}
\scriptsize
\setlength{\tabcolsep}{2.8pt}
\renewcommand{\arraystretch}{1.12}
\begin{tabular}{@{}p{0.14\columnwidth}p{0.24\columnwidth}cccc@{}}
\toprule
Family & Target / draft & NoSpec & TETRIS-Proxy & \system{} & Lift \\
\midrule
Qwen2.5 & 1.5B / 0.5B & 96.6$\pm$3.1 & 48.0$\pm$5.1 & 53.5$\pm$4.4 & \bad{$-$44.6\%} \\
Pythia~\cite{biderman2023pythia} & 2.8B / 160M & 96.3$\pm$2.2 & 146.5$\pm$2.9 & 148.6$\pm$4.4 & \pos{+54.3\%} \\
\bottomrule
\end{tabular}
\vspace{-0.5em}
\end{table}

Table~\ref{tab:model-pair-pressure} isolates pressure cases that would
be too detailed for the main result table. The KV-budget rows show that
the best eviction variant can change with both model pair and memory
budget, so these cells support the cache-regime boundary claim rather
than a universal eviction-rule ranking. The SLO-pressure rows report the
paired lift and confidence interval against NoSpec. The detector column
therefore records whether the warm-up evidence would enable or disable
speculation under that stressed setting.

% Inlined from ../paper_assets/tables/model_pair_pressure.tex.
\begin{table*}[t]
\centering
\caption{Additional $N=32$ model-pair pressure checks across three seeds. KV rows report SLO-aware goodput in tok/s ($\uparrow$). SLO rows report paired \system{} lift over NoSpec with 95\% confidence intervals.}
\label{tab:model-pair-pressure}
\scriptsize
\setlength{\tabcolsep}{2.2pt}
\renewcommand{\arraystretch}{1.12}
\begin{tabular}{@{}p{0.10\textwidth}p{0.11\textwidth}p{0.12\textwidth}p{0.18\textwidth}p{0.16\textwidth}p{0.11\textwidth}p{0.08\textwidth}@{}}
\toprule
Family & Stressor & Setting & \system{} result & Best / reference & Lift or gap & Evidence \\
\midrule
\multicolumn{7}{@{}l}{\textit{KV-pressure: best eviction variant can change with pair and budget}} \\
\cmidrule(l){1-7}
Qwen2.5 & KV budget & 256 blocks & 53.2$\pm$6.1 & SpecSched-H2O 56.0$\pm$0.8 & \bad{$-$5.0\%} & goodput \\
Qwen2.5 & KV budget & 512 blocks & 50.9$\pm$4.1 & SpecSched-Longest 58.7$\pm$3.5 & \bad{$-$13.3\%} & goodput \\
Pythia & KV budget & 256 blocks & 131.6$\pm$5.8 & SpecSched-H2O 135.7$\pm$2.8 & \bad{$-$3.0\%} & goodput \\
Pythia & KV budget & 512 blocks & 148.2$\pm$2.2 & SpecSched-H2O 152.1$\pm$1.5 & \bad{$-$2.6\%} & goodput \\
\midrule
\multicolumn{7}{@{}l}{\textit{SLO-pressure: detector decision follows the sign of paired CI}} \\
\cmidrule(l){1-7}
Qwen2.5 & SLO multiplier & $\beta=1.25$ & \bad{$-$53.8\%} & NoSpec & [$-$62.6\%, $-$45.0\%] & \dOFF{OFF} \\
Qwen2.5 & SLO multiplier & $\beta=2.0$ & \bad{$-$31.6\%} & NoSpec & [$-$44.3\%, $-$19.0\%] & \dOFF{OFF} \\
Pythia & SLO multiplier & $\beta=1.25$ & \pos{+90.5\%} & NoSpec & [+76.8\%, +104.2\%] & \dON{ON} \\
Pythia & SLO multiplier & $\beta=2.0$ & \pos{+51.0\%} & NoSpec & [+50.3\%, +51.6\%] & \dON{ON} \\
\bottomrule
\end{tabular}
\vspace{-0.5em}
\end{table*}

Table~\ref{tab:scheduler-overhead} is a CPU-side microbenchmark of the
policy computation only. Its role is to check scaling with $N$ and
$k_{\max}$ under a controlled local model pair. It is not an
end-to-end serving comparison. The measured wall-time fractions remain
below one percent, indicating that the reported goodput differences are
not dominated by scheduler bookkeeping overhead.

% Inlined from ../paper_assets/tables/scheduler_overhead.tex.
\begin{table}[t]
\centering
\caption{CPU-side scheduler overhead microbenchmark on the 0.5B/0.5B local model pair. Values are microseconds per scheduling epoch. Blue shading scales with overhead within this table.}
\label{tab:scheduler-overhead}
\small
\setlength{\tabcolsep}{3.2pt}
\renewcommand{\arraystretch}{1.12}
\begin{tabular}{rrrrr}
\toprule
$k_{\max}$ & $N$ & Mean ($\mu$s) $\downarrow$ & P95 policy ($\mu$s) $\downarrow$ & Wall (\%) $\downarrow$ \\
\midrule
3 & 8 & \heat{6}{128.8} & \heat{5}{54.8} & 0.15 \\
3 & 16 & \heat{11}{151.3} & \heat{8}{65.5} & 0.19 \\
3 & 32 & \heat{25}{220.0} & \heat{30}{148.1} & 0.20 \\
3 & 64 & \heat{40}{294.8} & \heat{40}{184.2} & 0.28 \\
\addlinespace
5 & 8 & \heat{5}{125.2} & \heat{8}{65.6} & 0.15 \\
5 & 16 & \heat{10}{146.7} & \heat{8}{65.7} & 0.18 \\
5 & 32 & \heat{26}{224.7} & \heat{32}{156.1} & 0.22 \\
5 & 64 & \heat{34}{263.0} & \heat{35}{165.3} & 0.25 \\
\addlinespace
7 & 8 & \heat{5}{123.2} & \heat{5}{56.3} & 0.14 \\
7 & 16 & \heat{10}{147.9} & \heat{12}{82.0} & 0.18 \\
7 & 32 & \heat{27}{232.2} & \heat{35}{167.5} & 0.23 \\
7 & 64 & \heat{32}{256.6} & \heat{33}{159.3} & 0.26 \\
\bottomrule
\end{tabular}
\vspace{-0.5em}
\end{table}

\section{Route and Remote-KV Stress Simulation}
\label{app:route-kv-sim}

This appendix reports the route/KV simulation referenced in
Section~\ref{sec:stress-failure}. The simulation exercises distributed
controls that the single-GPU artifact cannot execute directly. It uses
the measured \system{} rows as its local-route baseline: the closed-loop
$N$ sweep anchors local goodput and tail latency, and the three-seed
KV-pressure rows anchor 512- and 1024-block cache behavior. It then
varies four routing modes (local-only, random, KV-aware, and oracle),
route delay in $\{0,2,5,10\}$ ms, KV budget, and concurrency. The
modeled quantities are remote-KV hit ratio, migration rate, and routing
oscillation. This experiment is a stress diagnostic for the route and
cache variables in the main-text joint objective. It is not a native cluster
measurement.

\begin{figure*}[t]
    \centering
    \includegraphics[width=\textwidth]{figures_scientific/route_kv_ablation_scientific.pdf}
    \caption{Routing and remote-KV stress simulation at $N=64$, 512 KV
    blocks, and 5 ms route delay. The panels show that random routing
    loses goodput and increases cache churn, while KV-aware routing
    recovers throughput by trading remote hits against migration cost.}
    \label{fig:route-kv-ablation}
\end{figure*}

\begin{figure}[t]
    \centering
    \includegraphics[width=\columnwidth]{figures_scientific/route_kv_joint_stress_scientific.pdf}
    \caption{Joint stress surface for KV-aware routing at 5 ms route
    delay. Goodput drops when high concurrency combines with the smaller
    KV budget, showing the expected load--cache coupling.}
    \label{fig:route-kv-stress}
\end{figure}

Figure~\ref{fig:route-kv-ablation} and
Table~\ref{tab:route-kv-sim} separate route choice from local cache
pressure. At $N=64$ and 512 KV blocks, local-only routing gives
51.8 tok/s. Random routing falls to 43.3 tok/s and reaches an
oscillation rate of 0.52 because requests move faster than remote KV
reuse can amortize migration. KV-aware routing raises goodput to
55.2 tok/s by increasing the remote-hit ratio to 0.41 while keeping
oscillation lower at 0.33. The oracle upper bound reaches 61.5 tok/s,
showing that route selection has material value but is limited by
remote-delay-induced tail latency. KV utilization and eviction rates
make the resource mechanism visible: the route that improves goodput is
not merely the route with the most remote reuse, but the one that
balances reuse against migration and cache churn. Figure~\ref{fig:route-kv-stress}
shows the joint stress surface for KV-aware routing: increasing the KV
budget from 512 to 1024 blocks raises goodput in every load cell, while
the largest degradation appears at $N=64$ with the smaller cache. This
matches the formulation's prediction that route cost and cache
residency must be optimized together.

% Inlined from ../paper_assets/tables/route_kv_stress_sim.tex.
\begin{table}[t]
\centering
\caption{Routing and remote-KV stress simulation at $N=64$, 512 KV
blocks, and 5 ms route delay. Values are calibrated simulation means
over five deterministic replicates.}
\label{tab:route-kv-sim}
\scriptsize
\setlength{\tabcolsep}{3.2pt}
\renewcommand{\arraystretch}{1.12}
\begin{tabular}{@{}lrrrr@{}}
\toprule
Route & Goodput $\uparrow$ & P99 $\downarrow$ & Hit & Osc. \\
\midrule
local-only & 51.8 & 1386 & 0.00 & 0.00 \\
random & 43.3 & 1835 & 0.15 & 0.52 \\
kv-aware & 55.2 & 1621 & 0.41 & 0.33 \\
oracle & 61.5 & 1607 & 0.59 & 0.29 \\
\bottomrule
\end{tabular}
\vspace{-0.5em}
\end{table}

\section{Open-Loop Controller Diagnostics}
\label{app:open-loop-diag}

We tested the diagnostic variants described in the main-text guardrail
discussion. In a fixed-seed Poisson probe,
\textit{\system{}-RegimeAware} improves SLO-aware goodput at $N=16$
from 76.96~tok/s for FixedK-1 and 78.23~tok/s for \system{} to
78.43~tok/s, while reducing P99 TPOT from 1052.6 to
1030.9~ms/token relative to \system{}. At $N=32$, however, it drops
to 68.47~tok/s, below FixedK-1 (69.58~tok/s) and \system{}
(69.84~tok/s). The stronger \textit{\system{}-AIMD} controller
explicitly uses pending delay and multiplicatively caps open-loop
draft length. It reduces selected draft length to 1.0 and lowers P99
TPOT at $N=16$ from 1050.8 to 1039.2~ms/token, but goodput falls from
78.23 to 77.75~tok/s. At $N=32$, it further reduces goodput from
68.71 to 67.36~tok/s and does not improve P99 TPOT. These diagnostics
show that queue pressure and arrival regime should be explicit
scheduler inputs, but also that queue-age ordering and a global
draft-cap controller are too blunt.

We therefore ran a third fixed-seed probe with two stronger open-loop
controllers. \textit{\system{}-Backpressure} exposes ready-pending
backlog through admission throttling and inserts arrival pressure,
oldest pending delay, acceptance prior, and KV pressure into the
per-request marginal utility and draft cap. At $N=16$, Backpressure
obtains 82.74~tok/s versus 82.12~tok/s for FixedK-1, while reducing
P99 TPOT from 1101.4 to 1024.5~ms/token and lowering target forwards
per token from 0.563 to 0.539. At $N=32$, it obtains 78.46~tok/s
versus 77.95~tok/s for FixedK-1 and lowers target forwards per token
from 0.569 to 0.563, but P99 TPOT increases from 1522.7 to
1763.9~ms/token. The selected draft length remains close to one
($1.12$ at $N=16$ and $1.02$ at $N=32$), meaning the controller mainly
prevents harmful deep drafting and uses speculative depth only
opportunistically.

\textit{\system{}-PressureAIMD} adds a stateful AIMD controller over
the batch-level extra-draft budget and switches to admission-stable
batching under open-loop backlog. A recovery-gated fractional-credit
variant avoids the earlier failure mode where AIMD reduced target
forwards but increased KV evictions or verifier service time. Across
seeds 7, 42, and 123, this controller consistently keeps KV evictions
equal to FixedK-1 and releases only shallow extra drafts. It also
slightly reduces target forwards per token in every completed cell.
At $N=16$, PressureAIMD obtains
71.14$\pm$8.10~tok/s versus 68.02$\pm$10.62~tok/s for FixedK-1, and
raises mean SLO attainment from 0.859 to 0.891. The paired goodput lift
is 5.14\%, but its 95\% confidence interval is
$[-9.63\%, 19.92\%]$. At $N=32$, it obtains
71.81$\pm$3.27~tok/s versus 71.93$\pm$3.44~tok/s for FixedK-1, with
the same mean SLO attainment (0.896) and a paired lift of $-0.16\%$
with confidence interval $[-1.11\%, 0.79\%]$. The measurements support
the design direction of using arrival pressure and pending delay
directly in control and marginal utility. They also show that the
current controller is best interpreted as a diagnostic stabilization
mechanism: it prevents harmful deep drafting and lowers target-model
work, but it is not yet a validated replacement for the main
algorithmic claim.


% ============================================================
%              REFERENCES
% ============================================================
\bibliographystyle{IEEEtran}
\bibliography{references}
\end{document}
